% BEGIN CR28_CONSTITUTIVE_INVERSE
\section{Recuperación funcional, composición y momento de Catalán}
\label{cr28:sec:restitucion-constitutiva}

La respuesta de las incidencias \(90/120\) conserva más información que
un valor escalar aislado. En la misma notación de
\eqref{eq:three-four-completion}, escribimos
\[
 f(s)=\frac{1+s+s^2}{1+s+s^2+s^3},\qquad
 \mathcal D=s\,\frac{d}{ds},\qquad
 \mathcal V=f(T^{30}).
\]
Las marcas \(P_\pm\) mantienen la orientación de los canales. El
transporte angular es el antecedente de la respuesta; su reconstrucción
funcional recupera esta representación sin sustituir el archivo de memoria
del que procede.

\begin{theorem}[Inversión de la respuesta constitutiva]
\label{cr28:thm:inversion-constitutiva}
La función \(f\) es una biyección decreciente de \((0,1)\) sobre
\((3/4,1)\). Su inversa \(\psi\) asigna a \(r\) la única raíz
\(s\in(0,1)\) de
\begin{equation}
 r s^3+(r-1)(1+s+s^2)=0.
 \label{cr28:eq:cubic}
\end{equation}
Con las dos respuestas etiquetadas \(r_\pm\), quedan determinados
\begin{equation}
 s_\pm=\psi(r_\pm),\quad q_\pm=s_\pm^{1/30},\quad
 x=-\frac12\log(q_+q_-),\quad
 y=\frac12\log\frac{q_-}{q_+}.
 \label{cr28:eq:angular-recovery}
\end{equation}
Para el operador completo \(S=\mathcal V^2\), positivo con espectro
en \((9/16,1)\), el cálculo funcional proporciona
\begin{equation}
 T=\bigl[\psi(S^{1/2})\bigr]^{1/30},\qquad
 \det T=\det\bigl[\psi(S^{1/2})\bigr]^{1/30}.
 \label{cr28:eq:operator-recovery}
\end{equation}
\end{theorem}
\begin{proof}
La derivada negativa de \eqref{eq:vacuum-monotonic34} y los límites
\(f(0^+)=1,\ f(1^-)=3/4\) prueban la biyección. Multiplicar \(f(s)=r\)
por su denominador positivo da \eqref{cr28:eq:cubic}. La raíz positiva
de orden \(30\) recupera \(q_\pm\); la suma y la diferencia de
\(\log q_\pm=-x\mp y\) dan las dos coordenadas angulares. Por
positividad, \(S^{1/2}=\mathcal V\). Aplicar \(\psi\) a cada
proyector espectral recupera \(T^{30}\) y la raíz positiva da \(T\).
Tomar su determinante concluye.
\end{proof}

Esta recuperación del operador completo no se extiende a un determinante
aislado. En efecto,
\[
 T_1=\operatorname{diag}(1/5,1/3),\qquad
 T_2=\operatorname{diag}(1/6,2/5)
\]
tienen determinante \(1/15\), pero poseen espectros distintos en la misma
cámara contractiva. La inyectividad demostrada implica
\(f(T_1^{30})^2\ne f(T_2^{30})^2\). Ambos canales y sus marcas
conservan la recuperación; su reducción a un producto pierde información.
% END CR28_CONSTITUTIVE_INVERSE

% BEGIN CR28_TRANSPORT_COMPOSITION
% Recuperación expositiva de INTERRELACIONES_ESTRUCTURALES_VACIO.md, §§6--8.
% Insertar después de la inversión cúbica y antes de las coordenadas logarítmicas.
\subsection{Composición transportada de las respuestas constitutivas}
\label{cr28:comp-sec}

El transporte angular y las incidencias $90$ y $120$ proceden de la
misma estructura APP--TRIT--TPK: las hojas conservan las evaluaciones y
sus acarreos; el régimen trítico conserva la orientación; el calendario
TPK transporta fase y memoria hasta los canales angulares. La respuesta
constitutiva aplica a esos canales la función $f$ de
\eqref{eq:vacuum-monotonic34}. Su inversión permite expresar, en las coordenadas
de respuesta, la composición de transportes del mismo marco. Se conserva
así el orden constructivo: transporte, composición y evaluación espectral.
La operación obtenida actúa sobre esta representación del estado enriquecido;
su inversa recupera los canales, mientras la historia completa de memoria
permanece en la estructura de procedencia.

Sean $s=q^{30}$ y $\psi=f^{-1}$. El exponente $30$ es el máximo común
divisor de las incidencias $90$ y $120$. La inversión cúbica determina
$\psi:(3/4,1)\longrightarrow(0,1)$. Adjuntando la respuesta de frontera
$f(1)=3/4$, extendemos la biyección a
\[
 \psi:D\longrightarrow(0,1],\qquad D=[3/4,1).
\]
En esa frontera se utiliza la función racional $f$; el cociente
$(1-q^{90})/(1-q^{120})$ tiene allí una singularidad removible.

\begin{proposition}[Monoide de respuestas y coordenada aditiva]
\label{cr28:comp-monoid}
La operación
\begin{equation}
 r\odot t=f\bigl(\psi(r)\psi(t)\bigr),\qquad r,t\in D,
 \label{cr28:comp-law}
\end{equation}
hace de $D$ un monoide conmutativo y cancelativo con elemento neutro $3/4$.
La aplicación
\begin{equation}
 \ell(r)=-\frac1{30}\log\psi(r)
 \label{cr28:comp-character}
\end{equation}
es un isomorfismo de ese monoide con $(\mathbb R_{\geq0},+)$.
El intervalo $(3/4,1)$ es cerrado bajo $\odot$; el neutro pertenece a
la extensión de frontera. El único elemento invertible dentro de $D$
es el neutro.
\end{proposition}
\begin{proof}
El producto de dos elementos de $(0,1]$ pertenece al mismo intervalo y
\[
 \psi(r\odot t)=\psi(r)\psi(t).
\]
Para tres respuestas, ambos órdenes de asociación tienen imagen
$\psi(r)\psi(t)\psi(u)$; la inyectividad de $\psi$ prueba la
asociatividad. La conmutatividad procede del producto escalar. Si
$r\odot t=r\odot u$, la positividad de $\psi(r)$ permite cancelar y
obtener $\psi(t)=\psi(u)$, luego $t=u$. El parámetro $1$ corresponde a
$3/4$ y proporciona la identidad. El logaritmo transforma el producto en
suma, y su rango en $(0,1]$ prueba que $\ell$ es una biyección sobre
$\mathbb R_{\geq0}$. En particular, para la respuesta de un canal $q$,
$\ell(r)=-\log q$. Si ambos parámetros son estrictamente menores que
$1$, su producto también lo es. Finalmente, dos números de $(0,1]$
tienen producto $1$ únicamente si ambos son $1$.
\end{proof}

\begin{proposition}[Cancelación admisible y prolongación inversa]
\label{cr28:comp-inverse}
Para $r,u\in D$, la ecuación $r\odot t=u$ tiene solución $t\in D$
exactamente cuando $u\geq r$. Esa solución es única y viene dada por
\begin{equation}
 t=f\!\left(\frac{\psi(u)}{\psi(r)}\right).
 \label{cr28:comp-cancellation}
\end{equation}
La misma función $f$ es una biyección de $(0,\infty)$ sobre $(0,1)$.
Con esta extensión de $\psi$, la ley~\eqref{cr28:comp-law} transporta el
grupo multiplicativo de los parámetros positivos al intervalo $(0,1)$.
La inversión de grupo satisface
\begin{equation}
 r^{\ominus}=\psi(r)r,\qquad
 r\odot r^{\ominus}=\frac34.
 \label{cr28:comp-group-inverse}
\end{equation}
\end{proposition}
\begin{proof}
La ecuación equivale a
$\psi(t)=\psi(u)/\psi(r)$. Este cociente pertenece a $(0,1]$
exactamente cuando $\psi(u)\leq\psi(r)$, que por monotonía decreciente
equivale a $u\geq r$. La biyección prueba existencia y unicidad.
La derivada de \eqref{eq:vacuum-monotonic34} es negativa para todo $s>0$;
los límites de $f$ en $0$ y en infinito son $1$ y $0$, respectivamente.
Esto establece la biyección extendida. Multiplicando numerador y
denominador por $s^3$ se obtiene
\[
 f(s^{-1})=\frac{s^3+s^2+s}{s^3+s^2+s+1}=s f(s).
\]
La inversa multiplicativa de $s=\psi(r)$ tiene, por tanto, respuesta
$\psi(r)r$. Su composición con $r$ corresponde al parámetro $1$.
\end{proof}

La prolongación a $s>1$, o equivalentemente a $q>1$, es la extensión
algebraica del lector. El sector físico contractivo del capítulo conserva
su dominio declarado. La coordenada $\ell$ se prolonga a un isomorfismo
con $(\mathbb R,+)$ y cambia de signo bajo la inversión de grupo.

\subsubsection{Realización operatoria en un marco de hojas común}

\begin{proposition}[Composición con proyectores marcados]
\label{cr28:comp-operator}
Sean $\mathcal R=\mathcal R^*$, $\mathcal R^2=I$ y
$P_\pm=(I\pm\mathcal R)/2$. Para $j=a,b$, sean
\[
 T_j=q_{j,+}P_++q_{j,-}P_-,\qquad
 \mathcal V_j=f(T_j^{30}),\qquad 0<q_{j,\pm}<1.
\]
La composición $T_{a\circ b}=T_aT_b$ satisface
\begin{equation}
 \mathcal V_{a\circ b}
 =f\bigl(\psi(\mathcal V_a)\psi(\mathcal V_b)\bigr).
 \label{cr28:comp-operator-law}
\end{equation}
Sus dos respuestas son $r_{a,\pm}\odot r_{b,\pm}$. Si
$T_j=\exp[-(x_jI+y_j\mathcal R)]$ con $x_j>|y_j|$, entonces
\begin{equation}
 T_aT_b=
 \exp[-((x_a+x_b)I+(y_a+y_b)\mathcal R)],
 \label{cr28:comp-angular-addition}
\end{equation}
y $x_a+x_b>|y_a+y_b|$.
\end{proposition}
\begin{proof}
La ortogonalidad $P_+P_-=0$ da
\[
 T_aT_b=q_{a,+}q_{b,+}P_++q_{a,-}q_{b,-}P_-.
\]
En particular, $T_a$ y $T_b$ conmutan. El cálculo funcional en estos
proyectores proporciona
$\psi(\mathcal V_j)=T_j^{30}$ y
$(T_aT_b)^{30}=T_a^{30}T_b^{30}$; aplicar $f$ prueba
\eqref{cr28:comp-operator-law}. Para toda función $g$ definida en los
dos autovalores se tiene, más explícitamente,
\[
 P_\pm g(T_j)=g(q_{j,\pm})P_\pm.
\]
Como $q_{j,\pm}=\exp[-(x_j\pm y_j)]$, sus productos prueban
\eqref{cr28:comp-angular-addition}. La desigualdad triangular concluye
$x_a+x_b>|y_a|+|y_b|\geq|y_a+y_b|$.
\end{proof}

Las coordenadas recuperadas pueden escribirse directamente como
\[
 x=\frac{\ell(r_+)+\ell(r_-)}2,\qquad
 y=\frac{\ell(r_+)-\ell(r_-)}2.
\]
La composición suma estas parejas. El intercambio simultáneo de hojas
permuta los dos canales de cada transporte, actúa como
$(x,y)\mapsto(x,-y)$ y es un automorfismo de la ley componente a
componente. La inversión de ambos transportes en la extensión algebraica
actúa como $(x,y)\mapsto(-x,-y)$. Las dos involuciones conmutan y
conservan funciones distintas: una intercambia las marcas de orientación;
la otra invierte el transporte.

Para representar el intercambio mediante conjugación se utiliza un
operador $L$ que satisfaga $L\mathcal R L^{-1}=-\mathcal R$. En la
representación
\[
 \mathcal R=\begin{pmatrix}0&1\\1&0\end{pmatrix},\qquad
 L=\begin{pmatrix}1&0\\0&-1\end{pmatrix},
\]
la multiplicación directa da $LP_\pm L^{-1}=P_\mp$ y, por tanto,
$LT_jL^{-1}=q_{j,-}P_++q_{j,+}P_-$. La conjugación por
$\mathcal R$ conserva sus propios proyectores. Una traza que reúne los
canales tampoco reemplaza a los proyectores en esta identidad funcional:
la información de sus etiquetas interviene en la composición.

La hipótesis de marco común es sustantiva. Por ejemplo, los operadores
positivos definidos
\[
 A_0=\begin{pmatrix}1/2&0\\0&1/3\end{pmatrix},\qquad
 B_0=\begin{pmatrix}1/2&1/6\\1/6&1/2\end{pmatrix}
\]
tienen autovalores positivos, pero
\[
 A_0B_0=\begin{pmatrix}1/4&1/12\\1/18&1/6\end{pmatrix}
\]
carece de autoadjunción. La proposición especifica la composición de
transportes en los proyectores comunes; su dominio no comprende una mezcla
arbitraria de medios materiales. La realización de una concatenación de
rutas HMT conserva además sus datos de memoria: el cierre algebraico del
cono de parámetros describe la representación y no selecciona por sí solo
nuevas rutas primarias.

\subsubsection{Producto de canales y composición de transportes}

La velocidad normalizada de un estado sigue determinada por
\[
 \widehat c=(r_+r_-)^{-1}.
\]
Aquí el producto ordinario reúne las dos respuestas de ese mismo estado.
La ley $\odot$ evalúa otra operación: la composición de dos transportes
en cada hoja. El siguiente cálculo exacto distingue sus argumentos:
\begin{equation}
 f(1/2)=\frac{14}{15},\qquad
 \frac{14}{15}\odot\frac{14}{15}=f(1/4)=\frac{84}{85}
 \ne\frac{196}{225}=\left(\frac{14}{15}\right)^2.
 \label{cr28:comp-rational-example}
\end{equation}
Las fracciones de este ejemplo son parámetros algebraicos de comprobación.
La expresión de $\widehat c$ y las identidades constitutivas anteriores
permanecen inalteradas. El contenido añadido es una ley explícita para
transportar la composición y recuperar su coordenada angular aditiva.

% END CR28_TRANSPORT_COMPOSITION

% BEGIN CR28_CATALAN_RECOVERY
\subsection{Restitución funcional del momento orientado}
\label{cr28:sec:catalan-recovery}

La relación diferencial~\eqref{cr28:eq:catalan-problem} admite una
inversa que conserva la familia completa de momentos. Su dominio es
la función constitutiva producida por los sectores del ciclo, con la
orientación de~\eqref{eq:vacuum-angular-channels}. Las incidencias $90=3\cdot30$ y
$120=4\cdot30$ fijan esa función antes de evaluarla en los canales
$s_\pm=q_\pm^{30}$. La condición en el origen corresponde a la
desaparición del momento cuando se extingue su argumento contractivo.
Esta condición determina también las constantes de integración.

\begin{theorem}[Restitución integral y unicidad]
\label{cr28:thm:catalan-recovery}
Sea $f(s)=(1-s^3)/(1-s^4)$, para $0<s<1$, la respuesta
de~\eqref{eq:three-four-completion}, y sea $\mathcal D=s\,d/ds$.
Existe una única función $F\in C^2((0,1))$ que satisface
\begin{equation}
 \mathcal D^2F(s)=\frac{s(1+s)}{1+s+s^2}f(s),
 \qquad \lim_{s\downarrow0}F(s)=0.
 \label{cr28:eq:catalan-problem}
\end{equation}
Está dada por
\begin{align}
 F(s)&=\int_0^s\log\frac{s}{t}\,
          \frac{1+t}{1+t+t^2}f(t)\,dt
       =\int_0^s\frac{\log(s/t)}{1+t^2}\,dt
       =\mathcal C_2(s),
 \label{cr28:eq:catalan-integral}\\
 f(s)&=\frac{1+s+s^2}{s(1+s)}\mathcal D^2F(s).
 \label{cr28:eq:catalan-inverse}
\end{align}
La extensión continua a $s=1$ recupera $F(1)=G_{\rm Cat}$.
\end{theorem}
\begin{proof}
La factorización $1+t+t^2+t^3=(1+t)(1+t^2)$ da
\[
 \frac{1+t}{1+t+t^2}f(t)=\frac1{1+t^2}.
\]
El integrando de~\eqref{cr28:eq:catalan-integral} es
integrable en cero, porque
\[
 0\leq F(s)\leq\int_0^s\log(s/t)\,dt=s.
\]
Por tanto $F(s)\to0$. La diferenciación en el intervalo abierto,
primero en un intervalo truncado y después por convergencia dominada,
produce
\[
 \mathcal DF(s)=\int_0^s\frac{dt}{1+t^2},
 \qquad \mathcal D^2F(s)=\frac{s}{1+s^2}
     =\frac{s(1+s)}{1+s+s^2}f(s).
\]
En particular, $F$ es de clase $C^2$ y satisface la ecuación.
Si $F_1,F_2$ son dos soluciones, la diferencia $H=F_1-F_2$
satisface $\mathcal D^2H=0$. Mediante $z=\log s$ se obtiene
$d^2H(e^z)/dz^2=0$, luego $H(s)=a\log s+b$.
La existencia del límite nulo cuando $s\downarrow0$ obliga primero
a $a=0$ y después a $b=0$. La condición sobre $F$ determina, pues,
la solución y también el límite $\mathcal DF(s)\to0$.

Para $s<1$, el desarrollo geométrico de $(1+t^2)^{-1}$ y la
integrabilidad absoluta de sus términos ponderados permiten integrar:
\[
 \int_0^s t^{2k}\log(s/t)\,dt
     =\frac{s^{2k+1}}{(2k+1)^2},\qquad
 F(s)=\sum_{k\geq0}\frac{(-1)^ks^{2k+1}}{(2k+1)^2}.
\]
La serie representa la familia integral~\eqref{cr28:eq:catalan-integral}. La convergencia
absoluta y uniforme de esta última serie en $[0,1]$ autoriza el
paso al extremo $s=1$. También lo autoriza en la integral la cota
integrable $\log(1/t)$, extendiendo el integrando por cero para
$t>s$. Finalmente, despejar $f$ en la ecuación diferencial da
\eqref{cr28:eq:catalan-inverse}, con denominadores
positivos en $(0,1)$.
\end{proof}

El límite del momento y el límite de su derivada segunda se toman
en sus respectivas clases de convergencia. En particular,
\[
 \lim_{s\uparrow1}\mathcal D^2\mathcal C_2(s)=\frac12
\]
es el límite radial de Abel de la serie diferenciada. La evaluación
de $\mathcal C_2(1)$ utiliza directamente su serie absolutamente
convergente. Esta distinción conserva el mismo operador de profundidad
en el interior y en su lectura de frontera.

\subsection{Compatibilidad de la restitución con canales y velocidad}
\label{cr28:sec:channels-velocity}

El teorema restituye una función a partir de otra función previamente
construida. La inversión cúbica~\eqref{cr28:eq:cubic} actúa, en cambio,
sobre las evaluaciones $r_\pm$: recupera los argumentos $s_\pm$
dentro de esa familia fijada. La composición de ambas operaciones
conserva la orientación del cuarto de giro, las etiquetas $P_\pm$
y las incidencias del transporte. Así, $G_{\rm Cat}=\mathcal C_2(1)$
es la lectura extrema de un objeto funcional que participa también
en las respuestas interiores $r_\pm=f(s_\pm)$.

Para $\mathcal V_m=\mathcal V\otimes I_{9^m}$ y
$\tau_m=\operatorname{Tr}/(2\cdot9^m)$, la lectura de velocidad
mantiene exactamente la normalización de
\eqref{eq:vacuum-traces}:
\begin{equation}
 \exp[-2\tau_m(\log\mathcal V_m)]
   =\frac1{r_+r_-}
   =(\det\mathcal V)^{-1}=\widehat c.
 \label{cr28:eq:velocity-normalized}
\end{equation}
En efecto, cada autovalor $r_\pm$ se repite $9^m$ veces, de modo
que $2\tau_m(\log\mathcal V_m)=\log r_++\log r_-$.
La identidad utiliza el determinante de la realización de dos hojas
y la traza normalizada de su amplificación. Las secciones
dimensionales ya fijadas dan entonces
\[
 c_{\rm int}=c_{\rm vac}=\widehat c\,u_C,
 \qquad \ell_0=c_{\rm vac}t_0,
\]
con el mismo estado, la misma carta y la misma duración elemental
de~\eqref{cr28:eq:same-velocity}.
La restitución funcional conserva esta sección de velocidad al
componerla con la acción y la longitud; la multiplicidad espectral
queda absorbida por la normalización declarada.

% END CR28_CATALAN_RECOVERY

% BEGIN CR28_CONSTITUTIVE_RADIAL_LINK
\subsection{Una misma sección constitutiva en acción, longitud y gravedad}
\label{cr28:sec:constitutive-radial}
Las relaciones \eqref{eq:vacuum-sections} y la restitución anterior
mantienen una única sección de velocidad:
\begin{equation}
 c_{\rm int}=c_{\rm vac}=\widehat c\,u_C,\qquad
 \ell_0=c_{\rm int}t_0,\qquad
 c_{\rm int}=(\mu_{\rm vac}\varepsilon_{\rm vac})^{-1/2}.
 \label{cr28:eq:same-velocity}
\end{equation}
La última igualdad es una identidad de secciones en bases dimensionales
compatibles. La demostración radial del
\cref{g26:sec:rigidez-radial-es} construye primero
\(L=(5759/23040)\alpha^{16}L_*\) y después determina
\begin{equation}
 G_{\rm ret}=\frac{c_{\rm int}^3L^2}{\hbar_{\rm ret}}
 =\frac{L^2}
 {\hbar_{\rm ret}(\mu_{\rm vac}\varepsilon_{\rm vac})^{3/2}}.
 \label{cr28:eq:constitutive-gravity}
\end{equation}
La segunda expresión resulta de sustituir la identidad constitutiva en
la primera; no introduce una nueva calibración de ninguno de los dos
canales. En coordenadas, si \(c_{\rm int}=c_*\widehat c\), entonces
\(G_{\rm ret}=c_*^3\widehat c^{\,3}L^2/\hbar_{\rm ret}\).
El factor de propagación interviene una sola vez. La recuperación de
Catalán, el transporte compuesto y la respuesta radial conservan así
el mismo origen angular, sus marcas y su normalización.
% END CR28_CONSTITUTIVE_RADIAL_LINK

